\chapter{Desarrollo e Implementación}

\section{Configuración del Entorno de Ejecución (Docker)}
Para asegurar un entorno de desarrollo consistente y facilitar el despliegue del sistema en la infraestructura de la UNSa, se ha optado por la contenerización mediante \textbf{Docker} y \textbf{Docker Compose}. Esta metodología permite aislar la aplicación de las particularidades del sistema operativo anfitrión, empaquetando todas las dependencias necesarias.

La arquitectura de servicios definida en el archivo \texttt{docker-compose.yml} consta de tres contenedores principales interconectados:

\begin{itemize}
    \item \textbf{Servicio App (PHP-FPM):} Construido sobre una imagen de PHP 7.4/8.0, este contenedor procesa la lógica de la aplicación Laravel. Se incluyen extensiones necesarias para el correcto funcionamiento de los reportes PDF y la comunicación con MySQL.
    \item \textbf{Servicio Nginx:} Actúa como servidor web y proxy inverso. Está configurado para servir los archivos estáticos y redirigir las peticiones dinámicas al contenedor de la aplicación.
    \item \textbf{Servicio DB (MySQL 5.7):} Gestiona la persistencia de datos. El volumen de datos está mapeado externamente para garantizar que la información de los turnos no se pierda al reiniciar los contenedores.
\end{itemize}

El uso de redes virtuales (\texttt{turnos\_network}) garantiza una comunicación segura y privada entre los servicios, exponiendo únicamente el puerto 8083 para el acceso externo al sistema.

\section{Estructura general del proyecto}
El proyecto sigue el estándar de \textbf{Laravel 8.x}, organizando el código de manera que la lógica de negocio esté claramente separada de la presentación y los datos. Los directorios clave identificados son:

\begin{itemize}
    \item \texttt{app/Models}: Contiene las entidades Eloquent como \texttt{Turno.php}, \texttt{Dependencia.php} y \texttt{Feriado.php}. Aquí se definen las relaciones (1:N, N:M) y los \textit{Accessors} para formatear datos.
    \item \texttt{app/Http/Controllers}: Organiza la lógica de respuesta. Se dividen en subcarpetas para separar las acciones del \texttt{frontend} (reserva pública) de las de administración (\texttt{DependenciasController}, \texttt{TurnosController}).
    \item \texttt{database/migrations}: Almacena el historial de cambios en la estructura de la base de datos, permitiendo recrear el esquema académico en cualquier momento.
    \item \texttt{resources/views}: Contiene las plantillas \textbf{Blade}. Se han implementado \textit{Layouts} maestros para mantener la coherencia visual en todo el sitio.
\end{itemize}

\section{Implementación del Frontend}
El frontend del Sistema de Turnos Académicos se ha diseñado bajo una premisa de "Simplicidad y Eficiencia".

\subsection{Vistas Blade y AdminLTE}
Para el panel administrativo, se integra la plantilla \textbf{AdminLTE 4 (Beta)}. Esta elección permite disponer de un menú lateral colapsable, tablas dinámicas y una interfaz profesional para el personal del Observatorio. Se han personalizado los componentes de AdminLTE para adaptarlos a la identidad visual de la Facultad de Ciencias Exactas.

\subsection{Reactividad con Vue.js 2}
En las secciones donde la interacción es crítica, como la selección de horarios en tiempo real, se utilizan componentes de \textbf{Vue.js 2}. Estos componentes consumen una API interna de Laravel para:
\begin{itemize}
    \item Cargar dinámicamente los trámites disponibles al seleccionar una dependencia.
    \item Refrescar los horarios libres sin necesidad de recargar toda la página al cambiar la fecha en el calendario.
    \item Validar los datos del solicitante (DNI, Celular) antes de enviar el formulario.
\end{itemize}

\section{Lógica del Backend y Procesos Críticos}
El backend orquestado por Laravel maneja procesos que garantizan la integridad del sistema:

\begin{itemize}
    \item \textbf{Validación de Cupos:} En el controlador \texttt{TurnosDependenciasReservasController}, se implementa una lógica que verifica que la cantidad de personas (especialmente en turnos grupales) no exceda la capacidad de la dependencia antes de confirmar la reserva.
    \item \textbf{Gestión de Feriados:} Se integra un motor de búsqueda que cruza la fecha seleccionada con la tabla de feriados para bloquear días de asueto académico o mantenimiento del Observatorio.
    \item \textbf{Generación de Comprobantes:} Mediante el paquete \texttt{laravel-snappy}, se renderiza una vista Blade específica como PDF, incorporando un código de turno único para que el visitante lo presente en su llegada.
\end{itemize}
